We characterize when finite-window quantum extreme learning machines (QELMs)
can replace quantum reservoirs and when relevant information outside the window
prevents arbitrarily accurate emulation. The same information limit applies to
real-world temporal tasks whenever the required past information cannot be
recovered from the available input window.
We study reset-window QELMs, which reset for each prediction, and quantum
reservoir computers (QRCs), which carry a quantum state forward.
Under our assumptions, sufficiently expressive QELMs with growing windows
and exact expectations can approximate QRC outputs arbitrarily well precisely
when the missing past's influence fades uniformly. We quantify the unavoidable
window error and relate the required context to observable forgetting.
Persistent memory can prevent finite-window replacement. We construct a
simple quantum reservoir in which an old input can affect outputs
arbitrarily far into the future. Every fixed-window predictor therefore
faces an unavoidable worst-case error.
A QRC can solve the corresponding memory-recall task using exact
expectations; no reset-window QELM can match it uniformly over unbounded
delays, regardless of its capacity.
Analytical benchmarks and simulations examine three distinct limitations:
missing history, insufficient features, and finite measurements.
